\scanpage{129}
Cette inégalité peut être remplacée par une égalité,
car il existe un \unsure{op\textsuperscript{r}|op.} \unsure{part\textsuperscript{t}} \unsure{is.|iso.} $U$ tel que
$UA=|A|$, \deleted{\unsure{puis, une}} et d'autre part, \unsure{comme|même} il
résulte aussitôt de la décomposition \unsure{spectrale|spectral} de
$|A|$ et du th. de Hölder \unsure{classique|classiques} sur les
espaces $L^p$ et $L^{p'}$, un opérateur hermitien
positif $H\in L^{p'}$ tel que $\operatorname{Tr}(H|A|)=\|A\|_p$,
soit \deleted{\unsure{$|$}}$\operatorname{Tr}(HUA)$\deleted{\unsure{$|$}}$=\|A\|_p$. Si on pose
$B=UA$, on aura $\|B\|_{p'}\leq\|U\|\ \|H\|_{p'}$ (formule
(1)) donc $\|B\|_{p'}\leq1$, ce qui prouve notre assertion.
Par suite on a bien
\[
\|A\|_p=\sup_{\|B\|_{p'}\leq1}|\operatorname{Tr}AB|\tag*{(7)}
\]
ce qui prouve que $\|A\|_p$ est bien une
\unsure{semi norme|seminorme|semi-norme}, comme borne supérieure de
semi-normes $A\to|\operatorname{Tr}AB|$

Une telle démonstration n'est plus \unsure{pratiquable|praticable}
pour $p<1$\,; \deleted{et} nous démontrerons la formule (4)
par voie analytique, grace au

\uline{Lemme} 1 \quad Soit $(\alpha_i)$ une suite \deleted{\unsure{de puiss}} \deleted{décroissante} de
nombres $\geq0$, \deleted{\unsure{de puissance $p$ième sommables}}
\inserted{et soit} \deleted{\unsure{Soit|soit}} $0<p<1$. Posons $f(z)=\prod_{i=1}^{\infty}(1+\alpha_i z)$.
Alors on a
\[
\sum\alpha_i^p=p\frac{\sin\pi p}{\pi}\int_0^\infty\frac{\log f(r)}{r^{1+p}}\,dr\tag*{(8)}
\]
(les deux membres pouvant éventuellement
être infinis)
% archive sheet number: 130
